\documentclass[12pt]{article}
\usepackage[T1]{fontenc}
\usepackage[utf8]{inputenc}
\usepackage{lmodern}
\usepackage{textcomp}
\usepackage{enumitem}
\usepackage{geometry}
\geometry{margin=1in}
\begin{document}
\begin{center}
Answer Key - Economics - Undergraduate Level Assessment 3 \
\small (Answers are ordered exactly as in the question paper.)
\end{center}

\section*{Multiple Choice}
\begin{enumerate}[label=\arabic*.]
\item \textbf{Answer:} A
\\ \textit{Options:} A) the minimum of average total cost; B) the point where average variable cost is equal to marginal cost; C) the profit-maximizing quantity; D) the maximum of average variable cost; E) the point of equilibrium between supply and demand
\\ \textit{Note:} Marginal cost intersects average variable cost at its minimum point because when marginal cost is below average variable cost, it pulls the average down, and when it is above, it drives the average up, thus establishing that the intersection represents the lowest point of average variable cost.
\item \textbf{Answer:} A
\\ \textit{Options:} A) when price goes up supply goes up; B) price and supply have no correlation; C) when supply goes up, demand goes down; D) demand creates its own supply; E) when demand goes up, supply goes down
\\ \textit{Note:} Say's law posits that an increase in price incentivizes producers to increase supply, as higher prices typically signal greater potential profitability in the market.
\item \textbf{Answer:} D
\\ \textit{Options:} A) more effective in X because the policy will increase net exports.; B) more effective in X because the policy will decrease net exports.; C) equally effective in X and Y.; D) less effective in X because the policy will decrease net exports.
\\ \textit{Note:} Expansionary monetary policy in an open economy with flexible exchange rates tends to lower interest rates, which can lead to a depreciation of the currency, making exports cheaper and imports more expensive, but the overall increase in net exports may be offset by reduced domestic consumption and investment, rendering the policy less effective compared to a closed economy where such exchange rate fluctuations do not occur.
\item \textbf{Answer:} D
\\ \textit{Options:} A) saving remains constant.; B) consumption decreases by more than \$X.; C) saving increases by exactly \$X .; D) saving increases by less than \$X .; E) consumption decreases by less than \$X.
\\ \textit{Note:} This answer is correct because, according to the marginal propensity to save (MPS), individuals tend to spend a portion of any increase in disposable income, resulting in savings increasing by a lesser amount than the increase in income.
\item \textbf{Answer:} E
\\ \textit{Options:} A) The invention of automated systems that have made traditional farming obsolete; B) Emergence of widespread diseases that make farming untenable; C) Shift in consumer preferences towards synthetic foods; D) A global decrease in the demand for agricultural products; E) People seek the higher incomes of the city, shorter hours, and better social life in town
\\ \textit{Note:} The correct answer highlights that individuals are motivated to leave farming for urban areas due to the promise of higher incomes, reduced working hours, and enhanced social opportunities, reflecting a shift in economic priorities and lifestyle preferences.
\end{enumerate}

\section*{True / False}
\begin{enumerate}[label=\arabic*.]
\setcounter{enumi}{5}
\item \textbf{Answer:} True
\\ \textit{Options:} A) True; B) False
\\ \textit{Note:} The statement is true because the bill rate, calculated based on the discount from the face value of the 91-day bill, aligns with the yield derived from the coupon issue, indicating consistent financial metrics.
\item \textbf{Answer:} True
\\ \textit{Options:} A) True; B) False
\\ \textit{Note:} Real incomes can remain stable or increase during inflation if wage adjustments keep pace with or exceed the rate of inflation, thereby preserving or enhancing purchasing power.
\item \textbf{Answer:} True
\\ \textit{Options:} A) True; B) False
\\ \textit{Note:} A white noise process is characterized by a constant variance and a lack of autocorrelation at all lags other than zero, meaning that its values are uncorrelated and have a constant average variability.
\item \textbf{Answer:} True
\\ \textit{Options:} A) True; B) False
\\ \textit{Note:} A negative or contractionary supply shock, such as an increase in oil prices, raises production costs, leading to higher inflation and lower output, which shifts the Phillips curve to the right, indicating a trade-off between higher inflation and higher unemployment.
\item \textbf{Answer:} False
\\ \textit{Options:} A) True; B) False
\\ \textit{Note:} While increasing the federal funds rate is indeed a tool used by the Federal Reserve to manage monetary policy, it is not the only primary tool, as the Fed also employs open market operations and reserve requirements to influence the money supply and economic conditions.
\end{enumerate}

\section*{Long Form Response}
\begin{enumerate}[label=\arabic*.]
\setcounter{enumi}{10}
\item \textbf{Answer:} Adam Smith's concept of the division of labor significantly enhances production efficiency by breaking down work into specialized tasks, allowing workers to focus on specific functions. This specialization leads to increased productivity as workers become adept at their assigned tasks, reducing the time lost in switching between different types of work. Moreover, the division of labor fosters innovation, as individuals can develop and refine techniques within their specialized roles. The implications for economic systems are profound; as industries adopt this model, they can scale operations more effectively, leading to greater output and economic growth. Ultimately, the division of labor underscores the importance of specialization, which drives efficiency and productivity in modern economies.
\\ \textit{Note:} correct because it effectively articulates how the division of labor leads to increased production efficiency through specialization, skill development, reduced task-switching time, and the potential for innovation in economic systems.
\item \textbf{Answer:} In an economy where the velocity of money remains constant, the relationship between real Gross National Product (GNP), the quantity of money, and the price level can be understood through the equation of exchange, MV = PQ, where M is the money supply, V is the velocity of money, P is the price level, and Q is the quantity of goods and services produced (real GNP). If the price level drops by 10\%, assuming the velocity of money does not change, this implies that either the quantity of money has decreased or real GNP has increased, or a combination of both. A decrease in the price level can stimulate demand, potentially increasing real GNP if consumers respond positively to lower prices, thus incentivizing production. However, if the quantity of money also decreases significantly, it could lead to deflationary pressures, adversely affecting economic growth. Therefore, maintaining a balance between these factors is crucial for economic stability and growth.
\\ \textit{Note:} correct because it accurately applies the equation of exchange to illustrate how a decrease in the price level, with constant velocity, necessitates an adjustment in either the money supply or real GNP to maintain equilibrium in the economy.
\end{enumerate}

\vfill
\noindent\textit{End of Answer Key}
\end{document}